\documentclass[../../main.tex]{subfiles}
\begin{document}

\begin{center}
  {\fontsize{48pt}{58pt}\selectfont \textbf{1997}\normalsize\textbf{年度}}
\end{center}

\begin{table}[htb]
  \centering
  \caption{出題テーマと難易度}
  \begin{tabular}{|c|p{9cm}|c|}
    \hline
    \textbf{問題番号} & \textbf{テーマ} & \textbf{難易度} \\
    \hline
    第1問 & 平行移動した放物線の接触条件 & \\
    \hline
    第2問 & 整数列の部分和の条件 & \\
    \hline
    第3問 & 球面上の4点と長方形 & \\
    \hline
    第4問 & 二次関数を巡回する連立方程式 & \\
    \hline
    第5問 & 色付きカードを増やす試行の期待値 & \\
    \hline
    第6問 & 対数螺旋の弧長と面積 & \\
    \hline
  \end{tabular}
\end{table}

\begin{table}[htb]
  \centering
  \caption{解答}
  \renewcommand{\arraystretch}{1.4}
  \begin{tabular}{|c|c|p{9cm}|}
    \hline
    \textbf{問題番号} & \textbf{小問} & \textbf{答え} \\
    \hline
    第1問 & -- & $-1<a<0$ \\
    \hline
    第2問 & (1) & $s_n=n^2$（証明問題） \\
    \cline{2-3}
    & (2) & $a_k=2k-1$（$1\le k\le n$） \\
    \hline
    第3問 & (1)，(2) & 証明問題 \\
    \hline
    第4問 & (1) & 証明問題 \\
    \cline{2-3}
    & (2) & 相異なる実数解は$8$組 \\
    \hline
    第5問 & -- & $\dfrac{n}{2}+3$ \\
    \hline
    第6問 & (1) & $L=\sqrt{2}(1-e^{-\pi/2})$ \\
    \cline{2-3}
    & (2) & $S=\dfrac{1-e^{-\pi}}{4}$ \\
    \hline
  \end{tabular}
\end{table}

\paragraph{試験形態}
6問．試験時間は要確認．

\end{document}
